\section{Physical blocking and transfer idempotence}

\begin{definition}[Physical blocking isometry]
    \label{def:physical_blocking_isometry}
    \lean{MPSTensor.HasPhysicalBlockingIsometry}
    \leanok
    An MPS tensor $A$ is a \emph{renormalization fixed point} if there is an
    isometry $V\colon \C^d\to\C^{d^2}$, with coefficients
    $V_{(i_1,i_2),j}$ and $V^\dagger V=\Id$, such that, for all physical indices
    $i_1,i_2$,
    \begin{align}
        A^{i_1}A^{i_2}
        &=\sum_{j=0}^{d-1}V_{(i_1,i_2),j}A^j.
        \label{eq:corr_physical_blocking_isometry}
    \end{align}
    This is the relation $AA=A$ in
    \cite[Definition~3.2]{Cirac2016MPDO_arXiv}.
\end{definition}

\begin{definition}[Idempotent transfer-map criterion]\label{def:mps_transfer_idempotence}
    \lean{MPSTensor.IsTransferIdempotent}
    \leanok
    \uses{def:transfer_map}
    The transfer map of $A$ satisfies the idempotence criterion when
    $\E_A\circ\E_A=\E_A$.
\end{definition}

\begin{theorem}[A physical blocking isometry gives transfer idempotence]
    \label{thm:physical_blocking_implies_transfer_idempotence}
    \lean{MPSTensor.isTransferIdempotent_of_kraus_isometry}
    \leanok
    \uses{def:mps_transfer_idempotence}
    Suppose there exists an isometry
    $V\colon \C^d\to\C^{d^2}$, written as coefficients $V_{(i_1,i_2),j}$, such
    that
    \begin{align}
        A^{i_1}A^{i_2}
        &=\sum_{j=0}^{d-1}V_{(i_1,i_2),j}A^j
        \label{eq:corr_kraus_isometry_decomp}
    \end{align}
    for all $i_1,i_2\in\{0,\ldots,d-1\}$.
    Then $\E_A^2=\E_A$.
\end{theorem}
\begin{proof}\leanok
    Substituting the decomposition (\ref{eq:corr_kraus_isometry_decomp}) into
    the double Kraus sum gives, for every bond matrix $X$,
    \begin{align}
        \E_A^2(X)
        &=\sum_{i_1,i_2}(A^{i_1}A^{i_2})X
            (A^{i_1}A^{i_2})^\dagger
        \notag\\
        &=\sum_{j,k}\left(\sum_{i_1,i_2}
            V_{(i_1,i_2),j}\overline{V_{(i_1,i_2),k}}\right)A^jX(A^k)^\dagger
        \notag\\
        &=\sum_{j,k}(V^\dagger V)_{j,k}A^jX(A^k)^\dagger
          =\sum_{j,k}\delta_{j,k}A^jX(A^k)^\dagger
          =\sum_jA^jX(A^j)^\dagger
          =\E_A(X).
        \notag
    \end{align}
    Hence $\E_A^2=\E_A$.
\end{proof}

\begin{theorem}[Kraus-isometry criterion for transfer idempotence]%
    \label{thm:kraus_isometry_iff_transfer_idempotence}
    \lean{MPSTensor.isTransferIdempotent_iff_kraus_isometry}
    \leanok
    \uses{def:mps_transfer_idempotence,
        thm:physical_blocking_implies_transfer_idempotence}
    The transfer map of $A$ is idempotent if and only if there exists an
    isometry $V\colon \C^d\to\C^{d^2}$, written as coefficients
    $V_{(i_1,i_2),j}$ with $V^\dagger V=\Id$, such that
    \begin{align}
        A^{i_1}A^{i_2}
        &=\sum_{j=0}^{d-1}V_{(i_1,i_2),j}A^j
        \label{eq:corr_AA_isometry}
    \end{align}
    for all $i_1,i_2\in\{0,\ldots,d-1\}$.
\end{theorem}
\begin{proof}\leanok
    \uses{thm:physical_blocking_implies_transfer_idempotence,
        thm:kraus_rectangular_freedom'}
    The backward implication is
    Theorem~\ref{thm:physical_blocking_implies_transfer_idempotence}. For the
    forward implication, expanding $\E_A^2=\E_A$ as a Kraus sum gives, for
    every $X$,
    \begin{align}
        \sum_{i_1,i_2}(A^{i_1}A^{i_2})X(A^{i_1}A^{i_2})^\dagger
        &=\sum_jA^jX(A^j)^\dagger.
        \notag
    \end{align}
    Theorem~\ref{thm:kraus_rectangular_freedom'} then supplies the isometry $V$
    relating the two Kraus families.
\end{proof}

\begin{theorem}[Physical and transfer-map fixed-point criteria]
    \label{thm:physical_blocking_isometry_iff_transfer_idempotence}
    \lean{MPSTensor.isTransferIdempotent_iff_hasPhysicalBlockingIsometry}
    \leanok
    \uses{def:physical_blocking_isometry, def:mps_transfer_idempotence}
    The transfer map of an MPS tensor is idempotent if and only if the tensor
    has a physical blocking isometry.
\end{theorem}
\begin{proof}\leanok
    \uses{thm:kraus_isometry_iff_transfer_idempotence}
    This is Theorem~\ref{thm:kraus_isometry_iff_transfer_idempotence}, with the
    physical blocking relation written as
    Definition~\ref{def:physical_blocking_isometry}.
\end{proof}

\begin{theorem}[Limits of the renormalization flow]
    \label{thm:renormalization_flow}
    \uses{def:physical_blocking_isometry}
    A tensor appears as a limit of the two-site renormalization flow if and only
    if it has a physical blocking isometry. This is
    \cite[Theorem~3.1]{Cirac2016MPDO_arXiv}.
\end{theorem}

\begin{figure}[ht]
    \centering
    \begin{center}
        % A^{i_1}A^{i_2} = sum_j V^{i_1 i_2}_j A^j (eq:corr_kraus_isometry_decomp,
        % CPSV16 Thm 3.1, lines 396-412): two copies of A on physical legs
        % i_1, i_2 fuse through the isometry V into a single copy of A on
        % physical leg j.  V is a wide pill spanning both blocked columns;
        % up at={1,2} keeps one open up leg per column (i_1, i_2), while
        % down at=center gives the summed index j one port.  Both
        % sides carry the same boundary signature
        % (virtual-west=1 | virtual-east=1 | physical-up=2).
        \begin{tenkz}[physical=up]
            \tn[up=$i_1$]{A} & \tn[up=$i_2$]{A}
        \end{tenkz}
        $ = $
        \begin{tenkz}[rows={op:none, ket}]
            \tn[pill, wide=2, up at={1,2}, down at=center,
                up={$i_1$,$i_2$}]{V} & \\
            \tn[wide=2]{A^j} &
        \end{tenkz}
    \end{center}
        \begin{center}
        % sum_{i_1,i_2} (V^dagger)^j_{i_1 i_2} A^{i_1} A^{i_2} = A^j
        % (CPSV16 Thm 3.1, lines 396--412, source image II_UAAA.png):
        % V^dagger has one open upper port j and two lower ports, at the
        % blocked columns, contracted separately with A^{i_1}, A^{i_2}.
        % The two pairleg events certify that neither contraction is lost;
        % the boundary signature (virtual-west=1 | virtual-east=1 |
        % physical-up=1 | physical-down=0) certifies that no second upper
        % port is introduced.
        \begin{tenkz}[rows={op:none, ket}]
            \tn[pill, wide=2, up at=center, down at={1,2},
                up=$j$]{V^\dagger} & \\
            \tn{A} & \tn{A}
        \end{tenkz}
        $ = $
        \begin{tenkz}[physical=up]
            \tn[up=$j$]{A}
        \end{tenkz}
    \end{center}
    \begin{align}
        \sum_{i_1,i_2}(V^\dagger)^j_{i_1,i_2}A^{i_1}A^{i_2}
        &=A^j.
        \label{eq:corr_AA_isometry_adjoint}
    \end{align}
    \caption{The fusion identity of the renormalization fixed-point
        characterization: two contracted copies of $A$ equal a single copy of
        $A$ with the isometry $V$ merging the summed index $j$ into the
        physical pair $(i_1,i_2)$, as stated in
        (\ref{eq:corr_AA_isometry}); conversely, applying $V^\dagger$
        to the two physical legs of the contracted pair recovers the single
        tensor, as stated in
        (\ref{eq:corr_AA_isometry_adjoint}). The source draws the
        isometry as $U$
        \cite[Theorem~3.1, lines~396--412]{Cirac2016MPDO_arXiv}.}
    \label{fig:rfp_kraus_isometry}
\end{figure}
